\section{Related work and what is not new here}
\label{sec:related}

\paragraph{Automorphisms in Pollard rho.}
Walking on classes of an automorphism group of order $A$ saves
$\sqrt A$ in the expected number of steps: negation
\citep{wienerzuccherato1998}, Frobenius on Koblitz curves
\citep{wienerzuccherato1998,glv2000,glv2001}, and the equivalence-class
treatment of \citet{duursma1999}. The ECC2K-130 campaign
\citep{bailey2009,bos2010,bernstein2010} uses the $\langle\sigma,-1\rangle$
classes of order $262$ with an iteration $P\leftarrow P+\sigma^{j(P)}(P)$
whose branch $j$ depends on the normal-basis Hamming weight, and
normalizes a distinguished point ``by taking the lexicographically
smallest value among $x_R,x_R^2,\dots,x_R^{2^{130}}$ in normal-basis
representation'' \citep{bailey2009}. That is the canonical form of
Definition~\ref{def:sf} with $\prec$ the numeric order of the normal
word; the form is not new. What is new is the exact shortcut for
computing it (Section~\ref{sec:canonicalization}), the measurement of
what it costs relative to the group operation it sits next to
(Section~\ref{sec:bench-canon}), and the observation that in a walk
that canonicalizes every step it is the dominant cost
(Section~\ref{sec:bench-rho}).

\paragraph{Frobenius-invariant factor bases.}
\citet{ggmp2020} answer the question of how the Frobenius endomorphism
can speed up index calculus on subfield curves: with a factor base
closed under $\pi$ a relation can be ``rewrit[ten] in terms of a smaller
factor base consisting of representatives for each Frobenius orbit'',
which ``allows to reduce the number of linearly independent relations we
need to collect by a factor $n$'' and reduces the linear-algebra cost by
$n^2$; they construct such bases from Galois-invariant subspaces via
linearised polynomials and from the Couveignes--Lercier isogeny
construction, show how Frobenius-shifted bases give the full $m!$
symmetry-breaking saving without an $m^2$ linear-algebra penalty
(their Theorem~3.2), and conclude that these improvements ``do not lead
to index calculus algorithms faster than Pollard $\rho$ on instances
used in practice''. The idea descends from Galois-invariant factor
bases in finite fields \citep{jouxlercier2006}. Everything in
Section~\ref{sec:ic-recap} is theirs. They define the orbit
representative ``in an ad-hoc manner'' and do not price it; this paper
is about that price and about the layouts that make the rewritten
relations cheap to produce and verify. The sibling repository's
structural finding that a Galois-invariant \emph{subspace} base has a
quotient gain of at most $n$ and is a net loss of $57$--$96$ bits at
$n=131$ and $73$--$122$ bits at $n=163$ for $m\le5$, because the
$\pi$-stable subspaces of $\FF_{2^n}$ have dimensions
$\{0,1,n-1,n\}$ when $\ord_n(2)=n-1$, is the exact-spectrum
counterpart of their Lemma~4.1 and is why the end-to-end evidence here
stops at $n=83$.\art{crypto-autoresearcher:knowledge/findings/KN-FIND-b9a41d}

\paragraph{Symmetries in the decomposition.}
\citet{gg2014} reduce the factor base by torsion orbits (a factor $1/4$
and then $1/8$ on binary Edwards curves, with the matrix cost falling
by the square), break symmetry for a net $(m-1)!$ and conclude that rho
``is still much faster'' at reasonable sizes;
\citet{fghr2014,fhjrv2014} use the symmetric group and small-order
torsion to lower the degrees of the summation polynomials. The
$2$-torsion symmetrisation measured in the sibling repository
(the \texttt{sym} arms of the sibling repository's symmetrised
end-to-end study) is in that line and is orthogonal to the Frobenius phase: it
lowers columns and degrees, never the trial count.

\paragraph{Index calculus on binary curves.}
The summation-polynomial method is \citet{semaev2004}, made
subexponential in a range of extension degrees by \citet{gaudry2009}
and \citet{diem2011}; the binary-field case with Gr\"obner bases is
\citet{fppr2012,petitquisquater2012}, with the variable choices of
\citet{hpst2013,shantzteske2013} and the linearised-polynomial view of
\citet{huang2015}. None of these papers reports the cost of orbit
canonicalization or a phase-carrying data layout, which is the gap this
paper fills; all of them agree with the verdict of
Section~\ref{sec:bench-e2e} that the method does not reach rho at
cryptographic sizes.

\paragraph{Least rotation.}
Booth's and Duval's algorithms compute the least rotation of a string
in linear time with a failure function; Lemma~\ref{lem:zerogap} is the
binary special case in which the alphabet has two letters and the
comparison is numeric, and it trades the linear scan for a constant
number of whole-word rotations. We are not aware of it being stated
before, but it is elementary and we would not be surprised.

\paragraph{Measurement discipline.}
The unit, the classes and the isolation gate of
Section~\ref{sec:bench-method} are taken from the companion
measurement paper \citep{ecdlpconstant}, which this paper should be
read alongside: it supplies the matched rho and the floors; this paper
supplies the representation underneath both.
